\section{Связь интеграла Римана и Лебега}
\begin{theorem}
    Пусть $f \in \mathcal{R}([a,b])$.
    Тогда $f \in \widetilde{L}_1([a,b], \mathcal{L}^1)$.
\end{theorem} 
\begin{proof}
    \[
    \begin{cases}
        \exists \lim\limits_{l(T) \to 0} \overline{\overline{s}}_f(T) = J \in \R, \\
        \exists \lim\limits_{l(T) \to 0} \underline{\underline{S}}_f(T) = J \in \R,
    \end{cases}
    \]
    — по определению того, что $f$ интегрируема по Риману.
    Выберем зафиксированную последовательность разбиений $T_n$:
    \[T_n := \left\{a + \frac{k(b-a)}{n}\right\}_{k=0}^n: \
    a = x_0 < x_1 = a + \frac{1}{n} < \dots < x_n = b.\]
    $\underline{\underline{S}}$ и $\overline{\overline{s}}$ — верхние и нижние суммы Дарбу, то есть для $\Delta_i := [x_{i-1}, x_i]$:
    \[\overline{\overline{s}}_f(T_n) := \sum\limits_{i=1}^{n}m_i l(\Delta_i),\]
    \[\underline{\underline{S}}_f(T) := \sum\limits_{i=1}^{n}M_i l(\Delta_i),\]
    где $m_i = \inf\limits_{\Delta_i}f(x)$, $M_i = \sup\limits_{\Delta_i}f(x)$.

    Определим для любого натурального $n$:
    \[\underline{f_n}(x) := m_i, \text{ если } x \in \widetilde{\Delta}_i := [x_{i-1}, x_i),\]
    \[\overline{f_n}(x) := M_i, \text{ если } x \in \widetilde{\Delta}_i := [x_{i-1}, x_i).\]
    Будем считать, что $\underline{f_n}(b) = \overline{f_n}(b) = 0$.
    Тогда
    \[\int_{[a,b]} \underline{f_n}(x) d \mathcal{L}^1(x) = \overline{\overline{s}}_f(T_n) \to J, \ n \to \infty,\]
    \[\int_{[a,b]} \overline{f_n}(x) d \mathcal{L}^1(x) = \underline{\underline{S}}_f(T_n) \to J, \ n \to \infty.\]
    Заметим, что $\underline{f_{n+1}}(x) \geqslant \underline{f_n}(x)$ для любых натуральных $n$ и для любых $x \in [a,b]$.
    Определим
    $\underline{g_n}(x) = f_n(x) - f_1(x)$, $x \in [a,b]$ — уже неотрицательная функция.
    Поскольку $f_1 \in \widetilde{L}_1([a,b], \mathcal{L}^1)$, то, применив лемму Беппо-Леви, получим:
    \[\exists \lim\limits_{n\to \infty} \int_{[a,b]} \underline{g_n}(x) d \mathcal{L}^1(x) = \int_{[a,b]} \underline{g}(x) d \mathcal{L}^1(x),\]
    где
    \[\underline{g}(x) = \lim\limits_{n\to \infty} \underline{f_n}(x) - f_1(x),\]
    где $g$ — ограниченная функция, поскольку ограничена $f$.

    Пусть $\underline{f}(x) = \lim\limits_{n\to \infty} \underline{f_n}(x), \ x \in [a,b]$, тогда
    \[\exists \lim\limits_{n\to \infty} \int_{[a,b]} \underline{f_n}(x) d \mathcal{L}^1(x) = \int_{[a,b]} \underline{f}(x) d \mathcal{L}^1(x).\]
    Аналогично доказывается
    \[\exists \lim\limits_{n\to \infty} \int_{[a,b]} \overline{f_n}(x) d \mathcal{L}^1(x) = \int_{[a,b]} \overline{f}(x) d \mathcal{L}^1(x),\]
    где $\overline{f}(x) := \lim\limits_{n\to \infty} \overline{f_n}(x)$ $\forall x \in [a,b]$.
    Итого
    \[J = \int_{[a,b]} \underline{f}(x) d \mathcal{L}^1(x) = \int_{[a,b]} \overline{f}(x) d \mathcal{L}^1(x) \implies
    \int_{[a,b]} (\overline{f}(x) - \underline{f}(x))d \mathcal{L}^1(x) = 0,\]
    но $\overline{f}(x) \geqslant \underline{f}(x)$ $\forall x \in [a,b]$, поэтому
    $\overline{f}(x) = \underline{f}(x)$ для $\mathcal{L}^1$-почти всех $x \in [a,b]$.
    \begin{multline*}
        \underline{f}(x) \leqslant f(x) \leqslant \overline{f}(x) \ \forall x \in [a,b]
        \implies\\\implies
        \underline{f}(x) = f(x) = \overline{f}(x) \ \mathcal{L}^1 \text{-п.в.} x \in [a,b]
        \implies\\\implies
        J = \int_{[a,b]} f(x) d \mathcal{L}^1(x).
    \end{multline*} 
\end{proof} 

\begin{remark}
    Вообще говоря, в силу критерия Лебега интегрируемости по Риману, $f$ просто измерима, а не только измерима в широком смысле.
\end{remark}

\begin{example}
    \[\{r_k\}_{k=1}^\infty = \Q \cap [0,1].\]
    Пусть $\epsilon \in (0,1)$ — зафиксированный.
    Пусть открытое:
    \[G = \bigcup\limits_{k=1}^{\infty} \left(r_k - \frac{\epsilon}{4^k}, r_k + \frac{\epsilon}{4^k}\right).\]
    Обозначим компакт:
    \[[0,1] \setminus G = K.\]
    \[\mathcal{L}^1(G) \leqslant \sum\limits_{k=1}^{\infty} \frac{2 \epsilon}{4^k} = 2\epsilon \cdot \frac{1}{4} \cdot \frac{1}{1-\frac{1}{4}} = \frac{2\epsilon}{3}
    \implies
    \mathcal{L}^1(K) \geqslant \frac{\epsilon}{3}.\]
    Определим 
    \[f(x) := \chi_K(x) \implies f \in \widetilde{L}_1([0,1],\mathcal{L}^1) \text{ и } 
    \int_{[0,1]}f(x) d \mathcal{L}^1(x) = \mathcal{L}^1(K)\]
    $f$ имеет разрывы на множестве $K$ положительной меры, а множетсво точек разрыва будет иметь положительную меру Лебега, следовательно, по критерию Лебега не будет интегрируемости по Риману.
\end{example}